\section{Síntesis y ampliación}

\begin{enumerate}
    \item Pasar de un LM preentrenado a un asistente requiere tres pasos: instruction finetuning, modelado de recompensas y optimización de la política.
    \item \textbf{RLHF}: entrenar un modelo de recompensa y optimizar con PPO; es la técnica central de modelos como ChatGPT.
    \item \textbf{DPO}: simplifica la optimización de preferencias a un problema de aprendizaje supervisado, sin necesidad de infraestructura de RL.
    \item La regularización KL evita el reward hacking y es un componente esencial de la optimización de preferencias.
    \item El RL muestra un gran potencial en el razonamiento de los LLM (matemáticas, programación).
    \item Las comparaciones de preferencias son más confiables que las calificaciones absolutas; el modelo de Bradley-Terry es el marco estándar.
    \item La fase de despliegue requiere un evidence dashboard, una drift alarm y un ops playbook, que en conjunto forman un ciclo cerrado de monitoreo en tiempo real.
    \item Cada drift intervention se registra en el runbook y se vincula con el governance checklist, para asegurar el registro de auditoría y una recuperación rápida.
\end{enumerate}

\subsection{Tabla de síntesis}

\begin{table}[H]
\centering
\begin{tabular}{p{0.2\textwidth} p{0.35\textwidth} p{0.35\textwidth}}
\toprule
Tema & Aprendizaje central & Acción de ingeniería \\
\midrule
Instruction Finetuning & Aporta las capacidades básicas & Inicializar el modelo con instruction data a gran escala \\
RLHF + KL & Logra el alignment mediante un modelo de recompensa y la restricción KL & Mantener la reference policy para evitar el reward hacking \\
DPO & La optimización de preferencias puede prescindir del RL & Entrenar la política directamente y reducir la complejidad de la infraestructura \\
Pipeline Monitoring & Requiere timeline + checklist & Establecer un dashboard y un guardrail checklist \\
Evidence Matrix & Validación conjunta del alineamiento con varias señales & Estructurar benchmark, drift, red team y logs \\
Deployment Trust & Pre-mortem/post-mortem + multi-role governance & Establecer un rollback runbook, definir responsabilidades y alinearlas con la evidence matrix \\
\bottomrule
\end{tabular}
\caption{Arquitectura de RLHF y puntos clave de ejecución}
\end{table}

\subsection{Lecturas adicionales}
\begin{itemize}
    \item Ouyang et al., \textit{Training Language Models to Follow Instructions with Human Feedback} (InstructGPT, 2022).
    \item Rafailov et al., \textit{Direct Preference Optimization: Your Language Model is Secretly a Reward Model} (DPO, 2023).
    \item DeepSeek-AI, \textit{DeepSeek-R1: Incentivizing Reasoning Capability in LLMs via Reinforcement Learning} (2025).
    \item Chung et al., \textit{Scaling Instruction-Finetuned Language Models} (FLAN-T5, 2022).
\end{itemize}

\end{document}
